\lesson{3}{от 25 мар 2025 12:45}{Продолжение}

Типы МЛ

При обучении с учтелем -- функционал качества может определяться как средняя ошибка ответов. Предполагается, что искомый алгоритм должен его минимизировать.

\begin{definition}[Задача классификации]
  задача присвоения меток объектам.
\end{definition}

\begin{definition}[Задача регрессии]
  \emph{Задача регресии} -- это задача предсказания одной непрерывной случайной величины $\nu : \Omega \rightarrow Y$ на основе значений других случайных величин $\Xi_i : \Omega \rightarrow X_i$.
\end{definition}

\subsection{Предобработка данных}

Составляет более 80\% времени, затрачиваемого на большинство проектов по искусственному интелекту и машинному обучению.

Проекты AI, относящиеся к распознаванию объектов / изображений, автономного транспорту и аннотированию текста и изображений, являются наиболее распространенными задачами при разметке данных.

\subsection{Нормализация данных}

\begin{definition}[Нормализация]
  \emph{Нормализация} -- это способ организации данных. В нормализованной базе нет повторяющихся данных, с ней проще работать и можно менять её структуру для разных задач. В процессе нормализации данные преобразуют, чтобы они занимали меньше места, а поиск по элементам был быстрым и результативным.
\end{definition}

\begin{definition}[Третья нормальная форма (3NF)]
  2NF + каждый неключевой атрибут нетразитивно функционально зависит от первичного ключа (каждое неключевое поле должно содержать информацию о ключе, полном ключе и ни о чем, кроме ключа).
\end{definition}

\begin{definition}[Очистка данных]
  \emph{Очистка данных} -- это процесс исправления или удаления неправильных, поврежденных, дублированных данных в пределах набора данных. Этапы процесса очистки данных зависят от конкретного набора данных.
\end{definition}

\subsection{Ошибки. Оценки и проверка гипотез}

\begin{definition}[Ошибка первого рода]
  \emph{Ошибка первого рода} -- ситуация, когда отвергнута верная нулевая гипотеза.
\end{definition}

\begin{definition}[Ошибка второго рода]
  \emph{Ошибка второго рода} -- ситуация, когда принята неверная нулевая гипотеза.
\end{definition}
